\documentclass[10pt,a4paper]{amsart}
\usepackage[margin=19mm]{geometry}
\usepackage{amsmath,amssymb,mathtools}
\usepackage[scr=rsfs,cal=euler]{mathalfa}
\usepackage{xfrac}
\usepackage[unicode,hidelinks,bookmarks=false]{hyperref}
\newcommand{\ie}{i.e.\ }
\newcommand{\cf}{cf.\ }
\newcommand{\resp}{respectively\ }
\newcommand{\Subsection}{\subsection}
\providecommand{\nobreakdash}{\nobreak\hbox{-}}
\setlength{\emergencystretch}{2em}
\multlinegap0pt
\usepackage{amssymb}
\usepackage{bm}
\usepackage{mathtools}
\usepackage{dsfont}
\usepackage{tikz}
\usepackage[normalem]{ulem}
\usepackage{xcolor}
\newtheorem{theorem}{Theorem}
\newtheorem{lemma}{Lemma}
\renewcommand{\d}[1]{\ensuremath{\operatorname{d}\!{#1}}}
\title{Shape analysis in Schauder spaces of the energy of heat problems in perturbed annular domains}
\author{Luca Di Persio, Riccardo Molinarolo}
\date{Source: arXiv:2607.00803v2 (CC BY 4.0)}
\begin{document}
\maketitle

\noindent\emph{Source and licence.} Shape analysis in Schauder spaces of the energy of heat problems in perturbed annular domains. Version arXiv:2607.00803v2 (CC BY 4.0), distributed by arXiv under the Creative Commons Attribution 4.0 International licence, http://creativecommons.org/licenses/by/4.0/, as recorded in the arXiv licence metadata for this exact version and shown on the abstract page https://arxiv.org/abs/2607.00803v2. Full licence notice: \texttt{licenses/arxiv-2607.00803-CC-BY-4.0.txt}.\par\smallskip

\noindent\textit{Editorial scope.} Counted unit: the article's theorem dir (source0.tex lines 1262--1277). The article's complete original proof of it is delivered with it (source0.tex lines 1279--1281, one \texttt{proof} environment of 3 lines). No mathematics is written, repaired or re-derived: the statement and the proof are the article's own lines, in the article's own notation, and no retained line is substituted or re-flowed.  Retained uncounted as the source-local context the proof needs: lines 179--208; lines 198--206; lines 866--876; lines 901--911; lines 913--923; lines 929--948; lines 957--1001; lines 1070--1086; lines 1176--1186; lines 1204--1212; lines 1217--1244; lines 1248--1258.  Nothing was omitted from any retained span, so the package carries no omission record.  No retained line contains a citation, so the wrapper carries no bibliography and no bibliography entry.  No retained line contains a figure environment, an image-inclusion command or drawing code, so no image is delivered and the package declares no asset dependency.  Editorial formatting only, in addition: an AMS \texttt{amsart} preamble that restates the article's own declarations and macros used by the retained lines, namely the environments \texttt{theorem}, \texttt{lemma} on separate counters, together with the standard packages those lines need.  The retained environments print in this document as Theorem 1, Lemma 1 in order of appearance, whereas the article prints the same items with the numbering of its own full document.  The complete native source of the article as distributed in the arXiv e-print for this exact version is bundled unchanged as \texttt{sources/204189/source0.tex}.
\begin{theorem}\label{thm energy dependence abstract}
    Assume that there exist an open neighborhood $\widehat{\mathcal{W}}_0 \subseteq \mathcal{W}^\ast_0$ of $\phi_0$ in $\mathcal{A}^{\Omega^o}_{\partial\Omega^i}$ and a family of functions $\{ u_\phi \}_{\phi \in \widehat{\mathcal{W}}_0}$ such that the following holds:
    \begin{itemize}
        \item[(i)] for all $\phi \in \widehat{\mathcal{W}}_0$, $u_\phi \in C_{0}^{\frac{1+\alpha}{2}; 1+\alpha}([0,T] \times (\overline{\Omega^o} \setminus \Omega^i[\phi] ))$ and satisfy
        \begin{equation}\label{princeqpertu family}
            \partial_t u - \Delta u = 0  \quad \text{in } ]0,T] \times (\Omega^o \setminus \overline{\Omega^i[\phi]});
        \end{equation}

        \item[(ii)] the map from $\widehat{\mathcal{W}}_0$ to $C_{0}^{\frac{1+\alpha}{2}; 1+\alpha}([0,T] \times (\overline{\Omega^o} \setminus \Omega^i))$, which takes $\phi$ to the function $U_\phi$ given by 
        \begin{equation}\label{eq: U_phi := u_phi Psi phi}
        U_\phi(t,x):= \left( u_\phi \circ \Psi[\phi]^T \right) (t,x) \quad \text{for } (t,x) \in \overline{\Omega^o} \setminus \Omega^i,
        \end{equation}
        where $\Psi[\phi]^T(t,x):=(t,\Psi[\phi](x))$, is of class $C^\infty$.
    \end{itemize}
    Then, the map from $\widehat{\mathcal{W}}_0$ to $C^1([0,T])$, which takes $\phi$ to the function $e_\phi$ given by
    \begin{equation*}
    e_\phi(t) := \frac{1}{2} \int_{\Omega^o \setminus \overline{\Omega^i[\phi]}} (u_\phi(t,y))^2 \d y \quad \forall t \in [0,T],
    \end{equation*}
    is of class $C^\infty$. In particular, for every $t\in(0,T)$,
    \begin{equation}\label{eq: derivative of e general}
    \begin{aligned}
        \frac{\d{}}{\d t}e_\phi(t) &= - \int_{\Omega^o \setminus \overline{\Omega^i[\phi]}} |\nabla u_\phi(t,y)|^2 \d y + \int_{\partial\Omega^o} 
        \nabla u_\phi(t,y)\cdot \nu_{\Omega^o}(y) \, u_{\phi}(t,y) \d \sigma_y 
        \\
        & \quad
        - \int_{\partial\Omega^i} \nabla u_\phi(t,\phi(y))\cdot \nu_{\Omega^i[\phi]}(\phi(y)) \, u_\phi(t,\phi(y)) \, \tilde\sigma_n[\phi](y) \d\sigma_y,
    \end{aligned}
    \end{equation}
    and the right-hand side extends continuously to $[0,T]$.
\end{theorem}

    \begin{equation}
    \begin{aligned}
        \frac{\d{}}{\d t}e_\phi(t) &= - \int_{\Omega^o \setminus \overline{\Omega^i[\phi]}} |\nabla u_\phi(t,y)|^2 \d y + \int_{\partial\Omega^o} 
        \nabla u_\phi(t,y)\cdot \nu_{\Omega^o}(y) \, u_{\phi}(t,y) \d \sigma_y 
        \\
        & \quad
        - \int_{\partial\Omega^i} \nabla u_\phi(t,\phi(y))\cdot \nu_{\Omega^i[\phi]}(\phi(y)) \, u_\phi(t,\phi(y)) \, \tilde\sigma_n[\phi](y) \d\sigma_y,
    \end{aligned}
    \end{equation}

\begin{lemma}\label{lemma rappr}
    Let $\Omega$ be a bounded open connected subset of $\mathbb{R}^n$ of class $C^{1,\alpha}$, such that $\Omega^-$ is connected and $\overline{\Omega}\subseteq\Omega^o$. Then, the map from $C_0^{\frac{\alpha}{2}; \alpha}([0,T] \times \partial\Omega^o) \times C_0^{\frac{\alpha}{2}; \alpha}([0,T] \times \partial \Omega)$ to the space 
    \begin{equation*}
    \mathcal{X}_{\mathrm{heat}} := \left\{u\in C_{0}^{\frac{1+\alpha}{2}; 1+\alpha}([0,T] \times (\overline{\Omega^o} \setminus \Omega)) \colon \partial_t u  - \Delta u =0 \text{ in } ]0,T] \times \Omega^o \setminus \overline{\Omega} \right\}.
    \end{equation*}
    that takes a pair $(\xi,\gamma)$ to the function $u_{\Omega^o,\Omega}[\xi,\gamma]$ defined by
    \begin{equation}\label{eq: representation single layer}
    u_{\Omega^o,\Omega}[\xi,\gamma] := \left(v^+_{\Omega^o} [\xi] + v^-_{\Omega}[\gamma] \right)_{| [0,T] \times (\overline{\Omega^o} \setminus \Omega)},
    \end{equation}
    is bijective.
\end{lemma}

\begin{equation}\label{princeqpertu 2}
\begin{cases}
    \partial_t u - \Delta u = 0 & \quad\text{in } ]0,T] \times (\Omega^o \setminus \overline{\Omega^i[\phi]}), 
    \\
    \frac{\partial}{\partial\nu_{\Omega^o}} u(t,x) = g^o(t,x)& \quad \forall (t,x)\in [0,T] \times \partial \Omega^o, 
    \\
    \frac{\partial}{\partial\nu_{\Omega^i[\phi]}} u (t,x) = g^i(t,x,u(t,x)) & \quad \forall (t,x)\in  [0,T] \times \phi(\partial\Omega^i),
    \\
    u(0,\cdot)=0 & \quad \text{in } \overline{\Omega^o} \setminus \Omega^i[\phi],
\end{cases}
\end{equation}

\begin{equation}\label{princeq 2}
\begin{cases}
    \partial_t u - \Delta u = 0 & \quad\text{in } ]0,T] \times (\Omega^o \setminus \overline{\Omega^i}), 
    \\
    \frac{\partial}{\partial\nu_{\Omega^o}} u(t,x) = g^o(t,x)& \quad \forall (t,x)\in [0,T] \times \partial \Omega^o, 
    \\
    \frac{\partial}{\partial\nu_{\Omega^i}} u (t,x) = g^i(t,x,u(t,x)) & \quad \forall (t,x)\in  [0,T] \times \partial \Omega^i,
    \\
    u(0,\cdot)=0 & \quad \text{in } \overline{\Omega^o} \setminus \Omega^i.
    \end{cases}
\end{equation}

\begin{equation}\label{eq M Robin}
\begin{aligned}
\mathcal{M}_1[\phi,\xi,\gamma]
&:=\left(\frac12I+W_{\partial\Omega^o}^*\right)[\xi]
+\nu_{\Omega^o}\cdot\nabla v^-_{\Omega^i[\phi]}
\left[\gamma\circ(\phi^T)^{(-1)}\right]_{|[0,T]\times\partial\Omega^o}
-g^o
&&\text{on }[0,T]\times\partial\Omega^o,
\\
\mathcal{M}_2[\phi,\xi,\gamma]
&:=\left(-\frac12I+W_{\phi(\partial\Omega^i)}^*\right)
\left[\gamma\circ(\phi^T)^{(-1)}\right]\circ\phi^T
\\
&\quad +(\nu_{\Omega^i[\phi]}\circ\phi)\cdot
\left(\nabla v^+_{\Omega^o}[\xi]_{|[0,T]\times\phi(\partial\Omega^i)}\circ\phi^T\right)
\\
&\quad -\mathcal{N}_{g^i} \left( \phi, v^+_{\Omega^o}[\xi]_{|[0,T]\times\phi(\partial\Omega^i)}\circ\phi^T + V_{\phi(\partial\Omega^i)} \left[\gamma\circ(\phi^T)^{(-1)}\right]\circ\phi^T \right)
&&\text{on }[0,T]\times\partial\Omega^i,
\end{aligned}
\end{equation}

\begin{theorem}\label{sec 4: thm recall}
Let $\mathcal{M}$ be given by \eqref{eq M Robin}. Then the following holds.
\begin{itemize}
    \item[(i)] Let $(\phi, \xi,\gamma) \in \mathcal{A}^{\Omega^o}_{\partial\Omega^i} \times \mathcal{X}_0^{\frac{\alpha}{2}; \alpha}$. Then the function
    \begin{equation*}
        \mathbf{v}_{\Omega^o,\Omega^i[\phi]}\left[\xi,\gamma \circ (\phi^T)^{(-1)} \right] := \left( v^+_{\Omega^o} [\xi] + v^-_{\Omega^i[\phi]} \left[\gamma\circ (\phi^T)^{(-1)}\right] \right)_{| [0,T] \times (\overline{\Omega^o} \setminus {\Omega^i[\phi]})}
    \end{equation*}
    (cf.\, \eqref{eq: representation single layer}) is a solution in $C_{0}^{\frac{1+\alpha}{2}; 1+\alpha}([0,T] \times (\overline{\Omega^o} \setminus \Omega^i[\phi]))$ of problem \eqref{princeqpertu 2} if and only if 
    \begin{equation*}
        \mathcal{M}[\phi,\xi,\gamma] = (0,0).
    \end{equation*}
    In particular, there exists a unique pair $(\xi_0,\gamma_0) \in \mathcal{X}_0^{\frac{\alpha}{2}; \alpha}$ such that 
    \[
    \mathbf{v}_0 =  \mathbf{v}_{\Omega^o,\Omega^i[\phi_0]}\left[\xi_0,\gamma_0 \circ (\phi_0^T)^{(-1)} \right] = \left( v^+_{\Omega^o} [\xi_0] + v^-_{\Omega^i} [\gamma_0] \right)_{| [0,T] \times (\overline{\Omega^o} \setminus {\Omega^i})}
    \]
    is the unique solution in $C_{0}^{\frac{1+\alpha}{2}; 1+\alpha}([0,T] \times (\overline{\Omega^o} \setminus \Omega^i))$ of problem \eqref{princeq 2} and 
    \[
    \mathcal{M}[\phi_0,\xi_0,\gamma_0] = (0,0).
    \]

    \item[(ii)] The map $\mathcal{M}$ from $\mathcal{A}^{\Omega^o}_{\partial\Omega^i} \times \mathcal{X}_0^{\frac{\alpha}{2}; \alpha}$ to $\mathcal{X}_0^{\frac{\alpha}{2}; \alpha}$ is of class $C^{\infty}$.

    \item[(iii)] The partial differential of $\mathcal{M}$ with respect to $(\xi,\gamma)$ evaluated at the point $(\phi_0,\xi_0,\gamma_0)$, which we denote by $\partial_{(\xi,\gamma)}\mathcal{M}[\phi_0,\xi_0,\gamma_0]$ is a homeomorphism from $\mathcal{X}_0^{\frac{\alpha}{2}; \alpha}$ into itself.

    \item[(iv)] There exist two open neighborhoods $\mathcal{Q}_0$ of $\phi_0$ in $\mathcal{A}^{\Omega^o}_{\partial\Omega^i}$ and $\mathcal{H}_0$ of $(\xi_0,\gamma_0)$ in $\mathcal{X}_0^{\frac{\alpha}{2}; \alpha}$ and a $C^{\infty}$ map 
    \[
    \Xi := (\Xi_1,\Xi_2): \mathcal{Q}_0 \to \mathcal{H}_0
    \]
    such that the set of zeros of $\mathcal{M}$ in $\mathcal{Q}_0 \times \mathcal{H}_0$ coincides with the graph of the function $\Xi$. In particular,
    \begin{equation*}
    \mathcal{M}[\phi, \Xi_1[\phi],\Xi_2[\phi]] = 0  \quad \forall\phi \in \mathcal{Q}_0, \qquad  \Xi[\phi_0]= (\Xi_1[\phi_0],\Xi_2[\phi_0])= (\xi_0,\gamma_0).
    \end{equation*}

    \item[(v)] For each $\phi \in \mathcal{Q}_0$, the function $\mathbf{v}_\phi \in C_{0}^{\frac{1+\alpha}{2}; 1+\alpha}([0,T] \times (\overline{\Omega^o} \setminus \Omega^i[\phi] ))$ given by
    \begin{equation*}
    \begin{aligned}
        \mathbf{v}_{\phi} :=& \,\mathbf{v}_{\Omega^o,\Omega^i[\phi]}\left[\Xi_1[\phi],\Xi_2[\phi]\circ (\phi^T)^{(-1)}\right]
        \\
        =&\left( v^+_{\Omega^o} [\Xi_1[\phi]] + v^-_{\Omega^i[\phi]}
        \left[\Xi_2[\phi]\circ (\phi^T)^{(-1)}\right] \right)_{| [0,T] \times (\overline{\Omega^o} \setminus {\Omega^i[\phi]})}
    \end{aligned}
    \end{equation*}
    is a solution of \eqref{princeqpertu 2} and $\mathbf{v}_{\phi_0}= \mathbf{v}_0$.
    \end{itemize}
\end{theorem}

\begin{theorem}\label{thm smooth energy robin}
    Let $\mathcal{Q}_0$, $\Xi=(\Xi_1,\Xi_2)$ and $\mathbf{v}_\phi$ be as in Theorem \ref{sec 4: thm recall}. Then there exists an open neighborhood $\widehat{\mathcal W}_0$ of $\phi_0$ in $\mathcal A_{\partial\Omega^i}^{\Omega^o}$ such that the map from $\widehat{\mathcal{W}}_0$ to $C^1([0,T])$, which takes $\phi$ to $\mathbf{e}_\phi$ given by
    \begin{equation*}
    \mathbf{e}_\phi(t) := \frac{1}{2} \int_{\Omega^o \setminus \overline{\Omega^i[\phi]}} (\mathbf{v}_\phi(t,y))^2 \d y \quad \forall t \in [0,T]
    \end{equation*}
    is of class $C^\infty$. In particular, for every $t\in(0,T)$,
    \begin{equation}\label{eq: derivative of e 2}
    \begin{aligned}
        \frac{\d{}}{\d t}\mathbf{e}_\phi(t)
        &=-\int_{\Omega^o\setminus\overline{\Omega^i[\phi]}}|\nabla \mathbf{v}_\phi(t,y)|^2 \d y
        +\int_{\partial\Omega^o} \mathbf{v}_\phi(t,y)g^o(t,y) \d\sigma_y
        \\
        &\quad -\int_{\partial\Omega^i} \mathbf{v}_\phi(t,\phi(y)) \, g^i(t,\phi(y),\mathbf{v}_\phi(t,\phi(y))) \, \tilde\sigma_n[\phi](y) \d\sigma_y,
    \end{aligned}
    \end{equation}
    and the right-hand side extends continuously to $[0,T]$.
\end{theorem}

\begin{equation}\label{princeqpertu}
\begin{cases}
    \partial_t u - \Delta u = 0 & \quad\text{in } ]0,T] \times (\Omega^o \setminus \overline{\Omega^i[\phi]}), 
    \\
    u(t,x) = f^o(t,x)& \quad \forall (t,x)\in [0,T] \times \partial \Omega^o, 
    \\
    u(t,x) = f^i(t,\phi^{-1}(x)) & \quad \forall (t,x)\in  [0,T] \times \phi(\partial\Omega^i),
    \\
    u(0,\cdot)=0 & \quad \text{in } \overline{\Omega^o} \setminus \Omega^i[\phi],
    \end{cases}
\end{equation}

\begin{equation}\label{eq M Dirichlet}
\begin{aligned}
\mathcal{M}_1[\phi](\mu,\eta) & := V_{\partial\Omega^o}[\mu] + v^-_{\Omega^i[\phi]}
\left[\eta \circ (\phi^T)^{(-1)}\right]_{|[0,T] \times \partial\Omega^o} && \text{on } [0,T] \times \partial \Omega^o,
\\
\mathcal{M}_2[\phi](\mu,\eta) & :=
V_{\phi(\partial\Omega^i)}\left[\eta \circ (\phi^T)^{(-1)}\right] \circ \phi^T + v^+_{\Omega^o}[\mu]_{|[0,T] \times \phi(\partial\Omega^i)}\circ \phi^T &&\text{on } [0,T] \times \partial \Omega^i,
\end{aligned}
\end{equation}

\begin{theorem}\label{sec 5: thm recall dir}
Let $\mathcal{M}$ be given by \eqref{eq M Dirichlet}. Then the following holds.
\begin{itemize}
    \item[(i)] Let $(\phi,\mu, \eta) \in \mathcal{A}^{\Omega^o}_{\partial\Omega^i} \times \mathcal{X}_0^{\frac{\alpha}{2}; \alpha}$. Then, the function 
    \[
    u_{\Omega^o,\Omega^i[\phi]}\left[\mu,\eta\circ (\phi^T)^{(-1)}\right] := \left( v^+_{\Omega^o} [\mu] + v^-_{\Omega^i[\phi]} \left[\eta\circ (\phi^T)^{(-1)}\right] \right)_{| [0,T] \times (\overline{\Omega^o} \setminus {\Omega^i[\phi]})} 
    \]
    (cf. \eqref{eq: representation single layer}) is a solution in $C_{0}^{\frac{1+\alpha}{2}; 1+\alpha}([0,T] \times (\overline{\Omega^o} \setminus \Omega^i[\phi]))$ of problem \eqref{princeqpertu} if and only if 
    \begin{equation*}
    \mathcal{M}[\phi](\mu, \eta) = (f^o,f^i).
    \end{equation*}
    Moreover, for every $\phi \in \mathcal{A}^{\Omega^o}_{\partial\Omega^i}$, there exists a unique pair $(\Lambda_1[\phi],\Lambda_2[\phi]) \in \mathcal{X}_0^{\frac{\alpha}{2}; \alpha}$ such that 
    the function $\mathbf{u}_\phi \in C_{0}^{\frac{1+\alpha}{2}; 1+\alpha}([0,T] \times (\overline{\Omega^o} \setminus \Omega^i[\phi] ))$ given by
    \begin{equation*}
    \begin{split}
        \mathbf{u}_{\phi} :=& \,u_{\Omega^o,\Omega^i[\phi]}\left[\Lambda_1[\phi],\Lambda_2[\phi]\circ (\phi^T)^{(-1)}\right] 
        \\
        =&\left( v^+_{\Omega^o} [\Lambda_1[\phi]] + v^-_{\Omega^i[\phi]}
        \left[\Lambda_2[\phi]\circ (\phi^T)^{(-1)}\right] \right)_{| [0,T] \times (\overline{\Omega^o} \setminus {\Omega^i[\phi]})}
    \end{split}
    \end{equation*}
    is a solution of \eqref{princeqpertu}.

    \item[(ii)] The map $\mathcal{M}$ from $\mathcal{A}^{\Omega^o}_{\partial\Omega^i}$ to $\mathcal{L} \left( \mathcal{X}_0^{\frac{\alpha}{2}; \alpha}, \mathcal{X}_0^{\frac{1+\alpha}{2}; 1+\alpha} \right)$ is of class $C^{\infty}$. In particular, for every $\phi \in \mathcal{A}^{\Omega^o}_{\partial\Omega^i}$, $\mathcal{M}[\phi]$ is linear, continuous and invertible.

    \item[(iii)] The map $\Lambda$ from $\mathcal{A}^{\Omega^o}_{\partial\Omega^i}$ to $\mathcal{X}_0^{\frac{\alpha}{2}; \alpha}$ which takes $\phi$ to $\Lambda[\phi] = (\Lambda_1[\phi],\Lambda_2[\phi])$ is of class $C^{\infty}$.
    \end{itemize}
\end{theorem}

\begin{theorem}\label{sec 5: thm smootheness dir}
    Let $\mathbf{u}_\phi$ be as in Theorem \ref{sec 5: thm recall dir}. Let $\Omega_\mathtt{int}$ be a bounded open subset of $\Omega^o\setminus \overline{\Omega^i}$ of class $C^{1,\alpha}$ such that
    \[
    \overline{\Omega_\mathtt{int}} \subseteq  \overline{\Omega^o} \setminus \overline{\Omega^i}.
    \]
    Let $\mathcal{Q}_\mathtt{int} \subseteq \mathcal{A}^{\Omega^o}_{\partial\Omega^i}$ be an open neighborhood of $\phi_0$ such that  
    \begin{equation*}
        \overline{\Omega_\mathtt{int}} \subseteq \overline{\Omega^o} \setminus \overline{\Omega^i[\phi]} \quad \text{for all }\phi \in \mathcal{Q}_\mathtt{int}.
    \end{equation*}
    Then, the map from $\mathcal{Q}_\mathtt{int}$ to $C_{0}^{\frac{1+\alpha}{2}; 1+\alpha}([0,T] \times \overline{\Omega_\mathtt{int}})$ that takes $\phi$ to $\left(\mathbf{u}_{\phi}\right)_{|[0,T]\times\overline{\Omega_\mathtt{int}}}$ is of class $C^\infty$.
\end{theorem}

\begin{theorem}\label{thm smooth energy dir}
    Let $\mathbf{u}_\phi$ be as in Theorem \ref{sec 5: thm recall dir}. Then there exists an open neighborhood $\widehat{\mathcal W}_0$ of $\phi_0$ in $\mathcal A_{\partial\Omega^i}^{\Omega^o}$ such that the map from $\widehat{\mathcal{W}}_0$ to $C^1([0,T])$, which takes $\phi$ to $\mathbf{e}_\phi$ given by
    \begin{equation*}
    \mathbf{e}_\phi(t) := \frac{1}{2} \int_{\Omega^o \setminus \overline{\Omega^i[\phi]}} (\mathbf{u}_\phi(t,y))^2 \d y \quad \forall t \in [0,T]
    \end{equation*}
    is of class $C^\infty$. In particular, for every $t\in(0,T)$,
    \begin{equation}\label{eq: derivative of e dir}
    \begin{aligned}
        \frac{\d{}}{\d t}\mathbf{e}_\phi(t) &= -\int_{\Omega^o \setminus \overline{\Omega^i[\phi]}} |\nabla \mathbf{u}_\phi(t,y)|^2 \d y + \int_{\partial\Omega^o}
        \nabla \mathbf{u}_\phi(t,y)\cdot \nu_{\Omega^o}(y)\,f^o(t,y) \d\sigma_y
        \\
        &\quad -\int_{\partial\Omega^i} \nabla \mathbf{u}_\phi(t,\phi(y))\cdot \nu_{\Omega^i[\phi]}(\phi(y))\,f^i(t,y) \, \tilde\sigma_n[\phi](y) \d\sigma_y,
    \end{aligned}
    \end{equation}
    and the right-hand side extends continuously to $[0,T]$.
\end{theorem}

\begin{proof}
    In view of the same representation provided by Lemma \ref{lemma rappr} used for both problems \eqref{princeqpertu 2} and \eqref{princeqpertu}, the proof of Theorem \ref{thm smooth energy dir} is analogous to the proof of Theorem \ref{thm smooth energy robin}. In fact, it follows by the general Theorem \ref{thm energy dependence abstract} and by using the results contained in Theorems \ref{sec 5: thm recall dir} and \ref{sec 5: thm smootheness dir}. We just mention that \eqref{eq: derivative of e dir} follows by \eqref{eq: derivative of e general} using the boundary conditions satisfied by $\mathbf{u}_\phi$ (cf. \eqref{princeqpertu}).
\end{proof}

\end{document}
